\section{Problem Statement and Research Questions}
\label{sec:problem}

We study the trustworthiness of an individual repair attempt. Following the distinction
drawn in Section~\ref{sec:related}, we treat trustworthiness not as a developer's
subjective confidence but as a set of measurable properties of the artifact an agent
produces for a bug. This section makes those properties precise. We first ground them in a
concrete repair (Section~\ref{subsec:motivating}), then define them as predicates over a
(bug, repair attempt) pair (Section~\ref{subsec:axes}), characterize how they can come
apart (Section~\ref{subsec:dissociation}), and state the research questions that organize
the rest of the paper (Section~\ref{subsec:rqs}).

\subsection{A Motivating Example}
\label{subsec:motivating}

\begin{figure}[t]
\small
\textbf{(a) Gold patch: repair the source} (\texttt{sphinx/ext/autodoc/mock.py})
\begin{lstlisting}
 class _MockObject:
     __display_name__ = '_MockObject'
+    __name__ = ''
     def __init__(self, *args, **kwargs):
-        self.__qualname__ = ''
+        self.__qualname__ = self.__name__

 def _make_subclass(name, module, ...):
     attrs = {'__module__': module,
              '__display_name__': module + '.' + name,
+             '__name__': name,
              ...}
\end{lstlisting}
\textbf{(b) \texttt{minimax-m3} run 3: guard one consumer} (\texttt{sphinx/util/typing.py})
\begin{lstlisting}
 def _restify_py37(cls):
     elif hasattr(cls, '__qualname__'):
         if cls.__module__ == 'typing':
             return ':py:class:`~%s.%s`' % (cls.__module__, cls.__qualname__)
+        elif getattr(cls, '__sphinx_mock__', False):
+            return ':py:class:`%s`' % cls.__display_name__
         else:
             return ':py:class:`%s.%s`' % (cls.__module__, cls.__qualname__)
\end{lstlisting}
\caption{Two repairs for \texttt{sphinx-doc\_\_sphinx-9658}. Both pass the instance's test
suite. The gold patch (a) restores a mock object's name at its source, so every consumer
benefits. The agent's patch (b) guards a single downstream consumer, leaving the source and
the sibling renderer \texttt{\_restify\_py36} unfixed. Accuracy cannot tell them apart;
localization can.}
\label{fig:motivating}
\end{figure}

Consider \texttt{sphinx-doc\_\_sphinx-9658} from \sweBenchVerified{}. When Sphinx documents
a class whose base is a mocked import, it renders the base as \texttt{unknown.secret.}, a
dangling name with an empty final component, instead of \texttt{unknown.secret.Class}. The
cause lies in Sphinx's mock objects: \texttt{\_MockObject} carries no \texttt{\_\_name\_\_},
\texttt{\_make\_subclass} never sets one on the subclasses it builds, and \texttt{\_\_init\_\_}
hard-codes \texttt{self.\_\_qualname\_\_ = ''}. A mocked class therefore reaches the
cross-reference renderer, \texttt{restify}, with no usable name. Both annotators localize
the fault to its source in \texttt{mock.py}, and the developer's patch repairs the identity
there, so every consumer of a mock's name benefits at once (Figure~\ref{fig:motivating}a).

One agent run, \texttt{minimax-m3} run~3, resolves the bug by a different route
(Figure~\ref{fig:motivating}b). Rather than restoring the mock's name at its source, it adds
a special case to a single consumer: a branch in \texttt{\_restify\_py37} that detects a mock
and substitutes \texttt{\_\_display\_name\_\_}. The patch passes every test in the instance's
suite, and \sweBenchVerified{} records it as resolved, a verdict we confirmed by re-running
the suite on the patched repository. Judged by accuracy alone, the two repairs are
indistinguishable.

They are not. The agent's patch guards one reader of the mock's name and leaves the source
untouched, so the sibling renderer \texttt{\_restify\_py36}, and any other code that reads a
mock's \texttt{\_\_qualname\_\_}, still sees the empty identity; the gold patch fixes all of
them at once. The agent's localization reflects this: it points at \texttt{typing.py}, a
faithful downstream consumer that renders correct output for every genuine class and
mishandles mocks only because the name arriving there was already corrupted upstream.
Measured against the ground truth, the localization is a miss, and a deserved one. The
detail that makes this instance instructive is visible in the agent's trajectory: its own
reasoning names the true source (\texttt{\_MockObject.\_\_init\_\_} sets
\texttt{self.\_\_qualname\_\_} to the empty string) and then anchors the repair downstream
regardless. The artifact a developer would receive, a passing patch and its stated
localization, does not reflect where the fault lies, even though the agent had traced it.

This is the phenomenon \trustRepair{} is built to measure: resolution and localization can
diverge, and a test-passing bar cannot see the difference. The rest of this section defines
the axes that can.

\subsection{Trust Axes}
\label{subsec:axes}

\paragraph{Evaluation unit and ground truth.}
A \emph{bug} $t$ is an issue reported against a fixed base commit of a repository. For each
bug we assume a human-established ground truth consisting of a \emph{fault location}
$\loc^{\star}(t)$ (the code that must change to fix the bug) and a \emph{root-cause
rationale} $\rat^{\star}(t)$ (a natural-language account of why the bug occurs). How this
ground truth is constructed is deferred to Section~\ref{sec:benchmark}. A \emph{repair
attempt} $a$ for $t$ is the output of an agent run: a \emph{patch} $p$, a \emph{claimed
localization} $\loc$ (the code the attempt treats as faulty, as evidenced by its edits and
explanation), and a \emph{stated rationale} $\rat$ (the agent's account of the bug). We
define trustworthiness as three predicates over the pair $(t, a)$.

\paragraph{\accuracy{}.}
$\Acc(t, a) = 1$ iff the patch $p$ resolves the bug, that is, the repository with $p$
applied passes the instance's fail-to-pass and pass-to-pass test suites. Accuracy is the
criterion the field already uses; it is a necessary but, as we argue below, insufficient
condition for trust.

\paragraph{\interpretability{}.}
Interpretability asks whether the agent repairs the bug for the right reasons. It is the
conjunction of two predicates:
\emph{localization validity} $\Loc(t, a) = 1$ iff the claimed localization $\loc$ coincides
with the ground-truth fault location $\loc^{\star}(t)$, and
\emph{rationale alignment} $\Rat(t, a) = 1$ iff the stated rationale $\rat$ correctly
identifies the root cause $\rat^{\star}(t)$.
We write $\Int(t, a) = \Loc(t, a) \wedge \Rat(t, a)$. Rationale alignment is a judgement on
free-form text; we assess it on an ordinal rubric that we collapse to this binary predicate,
and we defer the judging procedure and its validation to Section~\ref{sec:method}.

\paragraph{Full trust.}
An attempt earns \emph{full trust} when it is accurate, correctly localized, and correctly
explained: $\FT(t, a) = \Acc(t, a) \wedge \Loc(t, a) \wedge \Rat(t, a)$. Full trust is the
property a developer would need in order to accept a patch \emph{because} it addresses the
actual defect, rather than merely because it passes the available tests.

\paragraph{\robustness{}.}
Let $\mathcal{T}$ be a set of semantics-preserving transformations: for $\tau \in
\mathcal{T}$, the variant $\tau(t)$ presents the same fault under an altered surface form,
and re-running the agent on $\tau(t)$ yields an attempt $a_\tau$. For a property $P \in
\{\Acc, \Loc, \Rat, \Int, \FT\}$, robustness is \emph{persistence} of the verdict under
transformation: $P$ is preserved on $(t, \tau)$ iff $P(\tau(t), a_\tau) = P(t, a)$. We
report robustness as the rate at which a property that holds on the original bug still
holds after transformation. The transformation set $\mathcal{T}$ is described in
Section~\ref{sec:benchmark} and the re-evaluation procedure in Section~\ref{sec:method}.

\subsection{Axis Dissociation}
\label{subsec:dissociation}

Because accuracy and interpretability are defined independently, an attempt can satisfy one
without the other. The dissociation central to this paper is \emph{fix-without-localizing},
$\Acc \wedge \lnot \Loc$ (which we also denote $\mathsf{fl\_wrong\_resolved}$): the patch
resolves the bug while its localization points at code in which the true fault does not lie,
as in Figure~\ref{fig:motivating}. A stronger form, \emph{fix-without-understanding}, adds a
failed rationale, $\Acc \wedge \lnot \Loc \wedge \lnot \Rat$: the patch resolves the bug
while the agent neither localizes nor explains the fault correctly. As we show in
Section~\ref{sec:results}, fix-without-localizing is common while its stronger form is rare,
because an agent that resolves a bug usually does describe its mechanism even when it
localizes elsewhere. The complementary case, $\mathsf{fl\_correct\_unresolved}$
($\Loc \wedge \lnot \Acc$), is one in which the agent identifies the faulty location but
fails to produce a passing fix.

These cases matter for different reasons. $\mathsf{fl\_wrong\_resolved}$ shows that a
passing patch is not evidence of correct localization, so accuracy overstates
trustworthiness. $\mathsf{fl\_correct\_unresolved}$ shows that correct
localization does not guarantee a fix, so localization carries information that accuracy
does not. Together they establish that the axes are not redundant restatements of one
another: each must be measured, and reported, in its own right. We study these dissociations
at the level of individual bugs rather than as an aggregate correlation, because a developer
confronts one repair at a time, and two agents with identical per-axis rates can still
disagree on which bugs they get right.

\subsection{Research Questions}
\label{subsec:rqs}

We organize the study around two research questions, each with two parts.

\begin{description}
  \item[\textbf{RQ1 (Trust axes and their alignment).}] \hfill
  \begin{itemize}
    \item \emph{RQ1a.} How do agents perform on each axis (accuracy, localization
    validity, and rationale alignment) and on full trust?
    \item \emph{RQ1b.} Do the axes agree at the level of individual bugs, or do agents
    produce fix-without-localizing ($\mathsf{fl\_wrong\_resolved}$) and the complementary
    $\mathsf{fl\_correct\_unresolved}$ repairs?
  \end{itemize}
  \item[\textbf{RQ2 (Robustness of trust).}] \hfill
  \begin{itemize}
    \item \emph{RQ2a.} Do accuracy and interpretability persist when a bug is subjected to
    semantics-preserving transformations?
    \item \emph{RQ2b.} Which axis degrades most under transformation?
  \end{itemize}
\end{description}

\noindent
\emph{Summary of Section~\ref{sec:problem}.} We defined trustworthiness as three
measurable axes over a (bug, repair attempt) pair: accuracy, interpretability (localization
validity and rationale alignment), and robustness. Together with the derived notion of full
trust and the two dissociations, $\mathsf{fl\_wrong\_resolved}$ and
$\mathsf{fl\_correct\_unresolved}$, these axes are what RQ1 and RQ2 investigate.
